\documentclass[11pt]{article}
\usepackage[margin=1in]{geometry}
\usepackage[T1]{fontenc}
\usepackage[utf8]{inputenc}
\usepackage{booktabs,hyperref,enumitem,xcolor}
\usepackage{amsmath,siunitx}

\newcommand{\repgithub}{https://github.com/raonialderete-eng/ems-mpc-datacenter}
\newcommand{\repzenodo}{10.5281/zenodo.22835283}

\title{Response to Reviewers\\[0.4em]
\large Predictive Energy Management of an Aggregated AI/HPC Data-Center Load with BESS}
\author{Raoni Alderete\\UFMS --- PPEE\\raoni.alderete@ufms.br}
\date{\today}

\begin{document}
\maketitle

\noindent\textit{Manuscript:} IEEE Access (decision: reject in present form with encouragement to resubmit).\\
\textit{Canonical TeX:} \texttt{04\_artigo/ieee\_access/artigo\_ems\_mpc\_datacenter\_access.tex}\\
\textit{Coverage source:} \texttt{PARECER\_ITEM\_A\_ITEM.md} (R1: 14 items; R2: 17; R3; R4.C1--C12; R5). This letter does \emph{not} reuse the old R1.15--R1.19 numbering.\\
\textit{Code:} \url{\repgithub}; archival DOI \url{https://doi.org/\repzenodo}.

\paragraph{How to read this letter.}
Each item is (a)~concern (paraphrase of the decision PDF, not a long quotation), (b)~response with evidence, (c)~action in the manuscript. Items that remain out of experimental scope are answered as explicit delimitations, not as ``Done''.

\section*{Associate Editor / process}

We thank the editor and the five reviewers. The resubmission package contains (i)~this letter, (ii)~a highlighted PDF, (iii)~clean PDF+TeX. Hardware claims are restricted as follows: SIL and UART HIL on a Terasic DE2-115 validate the \emph{embedded controller} on the same aggregated plant used in simulation. They are not a MW facility test, not UPS inner-loop validation, and not FIL (HDL Verifier). UART HIL of the heuristic is complete on four short scenarios. A 1200-step serial MPC HIL campaign and $T_s=\SI{1}{\second}$ ADMM on Cyclone~IV Lite are not claimed. A MATLAB row labelled \texttt{serial/mpc} at $90.98$\,kW / $57$\,s wall was a switch-capture error and is discarded.

%% =================================================================
\section*{Reviewer 1}

\subsection*{R1.1 Experimental / HIL validation}
\textbf{(a) Concern.} Lack of experimental or HIL evidence, especially for alleged zero unserved-power transients.

\textbf{(b) Response.} Desktop zero (filtered) remains a \emph{simulation} KPI on the aggregated plant. Additional evidence: SIL of the FPGA-ready copy ($N_p=12$, $N_c=4$, ADMM, same UART protocol) matches golden V5 on Group~A filtered peaks and sits $\approx 1$\,kW above the $350$\,kW bound on source loss (Table~\texttt{tab:sil}; \texttt{05\_resultados/hil\_de2115/}). UART HIL of the \emph{heuristic} on COM4 reproduced SIL on four short scenarios. ADMM UART executed tens of \texttt{RSP} frames until OpenCore Plus halted the time-limited Nios IP ($\approx 1$\,h); each QP exceeded $T_s$ on Nios soft-float. FIL was not implemented.

\textbf{(c) Action.} New Section~Embedded Validation in the manuscript. Abstract and limitations state SIL + heuristic UART HIL, and explicitly refuse FIL, completed MPC serial HIL, and 1\,s ADMM Lite.

\subsection*{R1.2 Measured AI/HPC traces}
\textbf{(a) Concern.} Use measured GPU/facility traces, not only synthetic steps.

\textbf{(b) Response.} Derived NVML (measured GPU power, 16 nodes, mapped to 600--950\,kW) and DIPLOEE-derived (simulated facility, resampled) windows are in \texttt{dataset\_v1/} with \texttt{manifest.json} hashes. They are not a feeder campaign owned by the authors, not synchronized GPU-utilization vs whole-facility power, and not 20 independent sites (windows overlap). Unit errata: metadata fields \texttt{raw\_min\_W}/\texttt{raw\_max\_W} for NVML are mislabelled after W$\to$kW; mapped traces used in tables are unaffected.

\textbf{(c) Action.} Trace paragraph in Scenarios; NVML/DIPLOEE rows in Table~\texttt{tab:punserved}.

\subsection*{R1.3 Timing sensitivity}
\textbf{(a) Concern.} Investigate timing errors; early alarms may beat the heuristic the wrong way.

\textbf{(b) Response.} The historical $\approx 699$\,kW at $-10$\,s is \emph{not} reproduced by the audited protocol. Consolidation \texttt{grupo\_timing.csv}: V5 at $-10$\,s is $350.00$\,kW (current sample remains measured; only the forecast shifts). None of 41 points on $[-20,+20]$\,s makes V5 worse than the heuristic $592.61$\,kW. $+10$\,s $\to 355.57$\,kW; $+20$\,s $\to 511.76$\,kW. An $N_{\mathrm{arm}}=8$ gate did not repair a Phase-1 $+10$\,s re-run and is \emph{not} sold as mitigation. The $699$\,kW change is a protocol/campaign difference, not an effect of the stochastic competitor.

\textbf{(c) Action.} Table~\texttt{tab:err\_time} and forecast-error subsection rewritten; abstract no longer claims that a 10\,s early alarm exceeds the reactive peak.

\subsection*{R1.4 Probabilistic timing / detection delay}
\textbf{(a) Concern.} Timing should be treated probabilistically, including realistic detection delays.

\textbf{(b) Response.} A 200-pair Monte Carlo draws forecast timing/magnitude errors on \emph{one} physical outage (Normal delays $\sigma=6$\,s clipped to $\pm 20$\,s). That law is not a detection-delay distribution. Separate detection-delay assays (1/5/10\,s) remain a distinct group. We do not equate the two.

\textbf{(c) Action.} Monte Carlo paragraph plus explicit delimitation in the information-policy subsection.

\subsection*{R1.5 Comparable tuning / heuristic variants}
\textbf{(a) Concern.} Comparable tuning methodology and alternative heuristic parametrizations.

\textbf{(b) Response.} Fairness campaign equalizes the common $\Delta u_{\max}$ knob (120/200/400\,kW/s). Equal \emph{names} do not equalize costs: myopic uses $w_{\Delta u}(u-u_{\mathrm{prev}})^2/P_{\mathrm{base}}^2$, V5 uses $w_{\Delta u}\|\Delta U\|^2$ without that division. Alternative heuristic charge-caps were not exhaustively re-optimized.

\textbf{(c) Action.} Fairness paragraph and heuristic subsection; closest horizon ablation is V5 with $N_c=N_p=1$.

\subsection*{R1.6 Weight sensitivity}
\textbf{(a) Concern.} Sensitivity of main weights and stability of conclusions.

\textbf{(b) Response.} Univariate $\times 0.5/1/2$ on five V5 weights, two scenarios, did not reverse Group~A ranking. It does not cover source-tracking, ramp, SOC, terminal and BESS-usage jointly. Large slacks are approximate priority, not lexicographic MPC.

\textbf{(c) Action.} Table~\texttt{tab:weights} plus the univariate caveat in the information-policy subsection.

\subsection*{R1.7 Converter/battery parameters vs hardware}
\textbf{(a) Concern.} Relate $\tau_b$ and efficiencies to hardware; discuss temperature, ageing, nonlinearities.

\textbf{(b) Response.} Sweep $\tau_b\in\{1,2,4\}$\,s keeps V5 on the case-study bound; ranking among five controllers is preserved. Varying $\tau_b$ does \emph{not} identify a cell or a UPS. Ageing/temperature are out of scope; $N_{\mathrm{EFC}}$ is a utilization proxy.

\textbf{(c) Action.} Table~\texttt{tab:tau} and model caveats; conclusion lists non-identified parameters.

\subsection*{R1.8 Multi-UPS / multi-bus / generators}
\textbf{(a) Concern.} Extension to multiple UPS units, converters, buses and generator coordination.

\textbf{(b) Response.} Not implemented. Qualitatively: per-unit bounds replace $750-400=350$\,kW; SOC sharing is absent from~\eqref{eq:dyn}; a genset could erase the $14.5833$\,kWh energy floor on a 150\,s outage.

\textbf{(c) Action.} Discussion-only paragraph in Related Work. No new architecture jobs.

\subsection*{R1.9 Partial criticality / shedding}
\textbf{(a) Concern.} Partial criticality, shedding priorities, redundancy.

\textbf{(b) Response.} The load is modelled as 100\% critical; there are no shedding classes. A class-split QP is not assembled.

\textbf{(c) Action.} Explicit hypothesis in model and related-work text; Appendix~A for $z$ (dimension 315).

\subsection*{R1.10 Computational headroom / real time}
\textbf{(a) Concern.} Do not claim real time from solver time alone; report cycle headroom.

\textbf{(b) Response.} On the i7-4510U campaign host, V5 load-band cycle p50/p95/max $\approx 0.022/0.031/0.242$\,s. Deterministic V5 recorded no $T_s$ overruns in principal short cases; the stochastic competitor recorded three first-sample overruns (not hidden). DE2-115 cycle statistics for 1200 ADMM steps are not a paper table. Warm-start is unused (empty $X_0$).

\textbf{(c) Action.} Computational subsection and Embedded Validation.

\subsection*{R1.11 Prolonged outage}
\textbf{(a) Concern.} Interpret long outages with autonomy, redundancy and generators.

\textbf{(b) Response.} Outage window $[600,7800)$\,s, $T_f=9600$\,s. Both controllers hit $\mathrm{SOC}_{\min}=20\%$. Unserved energy $1309.95$\,kWh (heuristic) vs $1312.33$\,kWh (V5): MPC leaves $\approx 2.38$\,kWh more unserved. Final SOC $\approx 24.75\%/24.73\%$ is after restoration/recharge, not the old $\approx 37\%$. Generator coordination is not simulated.

\textbf{(c) Action.} Table~\texttt{tab:outage} and prolonged-outage subsection.

\subsection*{R1.12 Positioning vs prior MPC/BESS}
\textbf{(a) Concern.} Differentiate the contribution from prior EMS-MPC studies.

\textbf{(b) Response.} Contribution is a controlled comparison protocol on an aggregated critical-load plant, not a new MPC theorem.

\textbf{(c) Action.} Table~\texttt{tab:lit} using already-cited works (Parisio, Wang 2023, Nair, Sepehrzad, Iqbal). No invented references.

\subsection*{R1.13 Reproducibility}
\textbf{(a) Concern.} Enough detail to reproduce loads, events and parameters.

\textbf{(b) Response.} Audited MATLAB, \texttt{jobs.json} seeds, \texttt{dataset\_v1} hashes, 11~Sep consolidation CSVs, SIL/HIL logs and FPGA-ready C are in the private GitHub tag.

\textbf{(c) Action.} Data Availability; URL \url{\repgithub}.

\subsection*{R1.14 Limit conclusions}
\textbf{(a) Concern.} Restrict claims to the model and to detectable/schedulable events.

\textbf{(b) Response.} Agreed.

\textbf{(c) Action.} Title, abstract and conclusions rewritten around an aggregated plant and detectable events.

%% =================================================================
\section*{Reviewer 2}

\subsection*{R2.1 Complete QP}
\textbf{(a) Concern.} Predictors, limits, slacks, recharge/event maps, normalizations.

\textbf{(b) Response.} Appendix~A transcribes the audited assembler: $z$ of dimension $315$, $P_{\mathrm{base}}=950$\,kW, $R_{\mathrm{base}}=100$, $w_u\|\Delta U\|^2$ \emph{not} divided by $P_b^2$, SOC linearization from the free response, zero extra SOC margin, no-preview fill = return-to-normal. \texttt{controle\_mpc\_qp\_v5.m} was not rewritten.

\textbf{(c) Action.} Appendix~A.

\subsection*{R2.2 Justify steps / restrict scope}
\textbf{(a) Concern.} Support synthetic steps with data or restrict title/scope.

\textbf{(b) Response.} Synthetic steps remain controlled isolators of $\tau_b$. Public traces complement them after affine mapping; they do not retro-calibrate the historical synthetics. Scope is restricted in the title.

\textbf{(c) Action.} Title + trace paragraph.

\subsection*{R2.3 Causal information classes}
\textbf{(a) Concern.} Separate scheduled load, planned interruption, alarm, and unannounced loss.

\textbf{(b) Response.} Current sample is measured; future $g$ of an unannounced outage is hidden. Detection delay is a separate assay. No-preview still fills future source-available / nominal limit (return-to-normal hypothesis). Stochastic early-event scenario without an announcement coincides with nominal.

\textbf{(c) Action.} Information-policy subsection.

\subsection*{R2.4 Four ablations / feedforward+feedback}
\textbf{(a) Concern.} Ablations and a feedforward+feedback competitor with the same forecasts.

\textbf{(b) Response.} Campaign includes foresight OFF, unit horizon, isolated event inequality. A fourth FF+FB structure was not added: Heur+FS, myopic, and V5 already isolate forecast vs local QP vs multi-step QP; unequal $\Delta u$ coefficients would remain.

\textbf{(c) Action.} Fairness/ablation paragraph; technical refusal of the extra competitor.

\subsection*{R2.5 Broad timing, missed events, repeated false alarms, gate}
\textbf{(a) Concern.} Wide timing, missed events, repeated false positives, combined errors, confirmation.

\textbf{(b) Response.} 41-point sweep, one false-alarm case, missed-event/unannounced loss, 200-pair MC. Repeated false alarms are \emph{not} in \texttt{criar\_jobs.py}. The arm-8 gate is not presented as a robust fix.

\textbf{(c) Action.} Forecast-error subsection; gate limitation explicit.

\subsection*{R2.6 Fallback / exitflags}
\textbf{(a) Concern.} Causes, instants, exitflags, fallback influence.

\textbf{(b) Response.} Fallback is part of the closed-loop policy and is tabulated (FB\%). Isolated solver-scale tests are not claimed to explain every historical \texttt{exitflag}.

\textbf{(c) Action.} Implementation paragraph + FB columns.

\subsection*{R2.7 SOC predictor inconsistency}
\textbf{(a) Concern.} Sign inconsistency frequency and SOC error vs piecewise model.

\textbf{(b) Response.} Accepted V5 solutions of seven principal scenarios: $100/8353$ steps ($1.20\%$) with a sign mismatch; worst error $0.044$\,pp; no bound violation of the piecewise reference.

\textbf{(c) Action.} Numbers in the foresight subsection.

\subsection*{R2.8 Myopic weights / active constraints}
\textbf{(a) Concern.} Myopic weight sensitivity, active constraints, command near steps.

\textbf{(b) Response.} $\Delta u_{\max}$ 120/200/400 on source-limit: $102.78$, $71.31$, $\approx 0$\,kW. Myopic-specific weight sweeps beyond that knob are not a full factorial. Coefficient mismatch is disclosed.

\textbf{(c) Action.} Fairness paragraph.

\subsection*{R2.9 Normalization / Hessian / hierarchical priority}
\textbf{(a) Concern.} Normalization, conditioning, hierarchical priority.

\textbf{(b) Response.} Reported gram indicator is $\mathrm{condest}(C^\top C+\varepsilon I)$ at the first step, not $\mathrm{cond}(C)$. Hessian $\approx 3.3\times 10^{5}$ on the main V5 load-band job. Lexicographic MPC is not implemented.

\textbf{(c) Action.} Information-policy and computational subsections.

\subsection*{R2.10 350 kW is case-specific}
\textbf{(a) Concern.} $750-400=350$\,kW is not a universal physical floor.

\textbf{(b) Response.} Agreed. It is a scenario-specific power bound. Over $150$\,s the persistent energy floor is $14.5833$\,kWh.

\textbf{(c) Action.} Terminology replaced throughout; ``physical floor'' removed from the TeX.

\subsection*{R2.11 UPS dynamics}
\textbf{(a) Concern.} Include UPS dynamics or restrict language to load continuity via BESS.

\textbf{(b) Response.} Language is restricted to load continuity on the aggregated model. DE2-115 with the same plant does not add electrical UPS fidelity.

\textbf{(c) Action.} Model, keywords, abstract, related work.

\subsection*{R2.12 Repeated/randomized evaluation}
\textbf{(a) Concern.} Randomize loads, SOC, converter, events, source limits and errors, with CI/quantiles.

\textbf{(b) Response.} 200 paired forecast errors and 40 mapped windows. The MC physical trajectory is the same in every draw. No joint randomization of event times and source limits. Not a population of data centers.

\textbf{(c) Action.} MC paragraph caveats.

\subsection*{R2.13 Cycle percentiles / warm-start / horizon scale}
\textbf{(a) Concern.} p50/p95/p99/max of cycle, assembly, iterations, scale with $N_p/N_c$, warm-start.

\textbf{(b) Response.} i7 cycle percentiles reported. Audited \texttt{quadprog} passes empty $X_0$; a warm-start trial matched iteration counts. Horizon-scale curves mixing another solver/board are not compared as a single scaling law.

\textbf{(c) Action.} Computational subsection: no warm-start claim.

\subsection*{R2.14 Long-outage final SOC}
\textbf{(a) Concern.} Explain final SOC above minimum after a long outage; give intervals.

\textbf{(b) Response.} Intervals $[600,7800)$\,s outage, $T_f=9600$\,s. SOC$_f$ $24.75\%/24.73\%$ after source return and recharge. Pre-event SOC $60\%/59.50\%$. A $7200$\,s ``recovery'' column is outage duration, not autonomous BESS restoration.

\textbf{(c) Action.} Table~\texttt{tab:outage}.

\subsection*{R2.15 Extra metrics}
\textbf{(a) Concern.} Times above thresholds, pre-event SOC, preparation energy, violations, recovery, throughput/EFC, reversals.

\textbf{(b) Response.} Basic metrics are in the consolidation CSVs. Preparation energy in an explicit window and reversal counts are \emph{not} consolidated; $u$ total variation is not a reversal count. $T_{\mathrm{gt1}}$ uses $\ge 1$; other thresholds use $>$. EFC is utilization, not lifetime.

\textbf{(c) Action.} Metrics section; remaining gaps named rather than filled with new jobs.

\subsection*{R2.16 Retracted / recent references}
\textbf{(a) Concern.} Replace retracted reference; check recent metadata.

\textbf{(b) Response.} The canonical TeX has no \texttt{sheng2026power} entry. Remaining bibliography uses already-verified keys; an online editorial-status audit of every 2025--2026 item is still the author's pre-upload check.

\textbf{(c) Action.} Bibliography retained without the retracted key.

\subsection*{R2.17 Permanent repository}
\textbf{(a) Concern.} Code, parameters, data, forecasts, seeds and results in a permanent archive.

\textbf{(b) Response.} Versioned code and derived data: \url{\repgithub}. Permanent archive: \url{https://doi.org/\repzenodo}. Raw $\approx 1$\,GB traces are not uploaded; \texttt{manifest.json} hashes are.

\textbf{(c) Action.} Data Availability and README of the tag.

%% =================================================================
\section*{Reviewer 3}

\subsection*{R3.1 Comparative table vs cited MPC works}
\textbf{(a) Concern.} Table versus prior works (information, plant, baselines, HIL/SIL, errors).

\textbf{(b) Response.} Qualitative positioning only; not a meta-analysis.

\textbf{(c) Action.} Table~\texttt{tab:lit} (Parisio, Wang 2023, Nair, Sepehrzad, Iqbal, this paper).

\subsection*{R3.2 Statistical support / CI}
\textbf{(a) Concern.} Confidence intervals on forecast tests.

\textbf{(b) Response.} 200 pairs; mean peak difference stochastic$-$V5 $= -9.04$\,kW, bootstrap 95\% CI $[-14.02,-4.79]$; seed $20260911$; fixed physical trajectory.

\textbf{(c) Action.} Monte Carlo paragraph.

\subsection*{R3.3 Justify weights}
\textbf{(a) Concern.} Justify weights and show sensitivity.

\textbf{(b) Response.} Load/event slacks dominate after $P_{\mathrm{base}}$ normalization (Table~\texttt{tab:weights}). Univariate sweep only; no universal stability claim.

\textbf{(c) Action.} Weights table + infopol caveat.

\subsection*{R3.4 First-order simplification vs ranking}
\textbf{(a) Concern.} Does the first-order converter change controller ranking?

\textbf{(b) Response.} Ranking is preserved under $\tau_b=1/2/4$\,s. That does not validate a richer converter/UPS model.

\textbf{(c) Action.} Table~\texttt{tab:tau} and explicit non-validation sentence.

\subsection*{R3.5 Identical 90.98 kW peaks}
\textbf{(a) Concern.} Equal heuristic peaks on load-bands and source-limit look like a copy-paste error.

\textbf{(b) Response.} Both expose a $150$\,kW first-step deficit through $\tau_b=2$\,s: $150\,e^{-1/2}=90.979599$\,kW, at $400$\,s vs $200$\,s. Scenarios are otherwise different.

\textbf{(c) Action.} Sentence in the metrics subsection.

%% =================================================================
\section*{Reviewer 4}

\subsection*{R4.C1 Scientific novelty}
\textbf{(a) Concern.} Novelty beyond implementing existing methods.

\textbf{(b) Response.} Applied, reproducible comparison protocol. The historical $-10$\,s / $699$\,kW counter-example is \emph{not} claimed as a result of this campaign.

\textbf{(c) Action.} Related-work table and timing rewrite.

\subsection*{R4.C2 Validation beyond simulation}
\textbf{(a) Concern.} HIL or operational measurements.

\textbf{(b) Response.} Same as R1.1: heuristic UART HIL closed; MPC serial campaign not closed; FIL not done. Measured NVML is an \emph{input} profile, not electrical validation of the simulated plant.

\textbf{(c) Action.} Embedded Validation.

\subsection*{R4.C3 Public GPU/HPC dataset}
\textbf{(a) Concern.} Use GPU/HPC measurements or an established public dataset.

\textbf{(b) Response.} Same as R1.2. Mapped 600--950\,kW is not a measured 600--950\,kW feeder.

\textbf{(c) Action.} Trace paragraph.

\subsection*{R4.C4 Probabilistic forecast/timing}
\textbf{(a) Concern.} Treat prediction/timing probabilistically.

\textbf{(b) Response.} Sweep + MC + three-scenario competitor. Assumed error law disclosed; $-10$\,s $\neq 699$\,kW in the audited grid.

\textbf{(c) Action.} Forecast-error subsection.

\subsection*{R4.C5 Established predictive EMS (e.g.\ stochastic MPC)}
\textbf{(a) Concern.} Add an established predictive EMS, e.g.\ stochastic MPC.

\textbf{(b) Response.} Fifth competitor: three scenarios, $\pi=1/3$, shared first command. It is this project's implementation of a method class, not a reproduction of a named published benchmark.

\textbf{(c) Action.} Stochastic subsection + Appendix scenario maps.

\subsection*{R4.C6 Weights of several orders}
\textbf{(a) Concern.} Justify widely different weights and robustness to tuning.

\textbf{(b) Response.} Approximate hierarchy via normalized slacks. Large weights are not strict priority. Univariate sweep only.

\textbf{(c) Action.} Table~\texttt{tab:weights}, infopol.

\subsection*{R4.C7 Real-time on representative hardware}
\textbf{(a) Concern.} Real-time execution evidence on representative hardware.

\textbf{(b) Response.} i7 cycle times $\neq$ DE2-115. Heuristic meets $T_s=1$\,s on the board in the four short cases. ADMM on Nios is demonstrated, not a 1\,s product. Stochastic MPC is not ported.

\textbf{(c) Action.} Embedded Validation; three stochastic overruns on the PC are disclosed.

\subsection*{R4.C8 Battery degradation / replacement cost}
\textbf{(a) Concern.} Incorporate degradation and replacement cost.

\textbf{(b) Response.} Out of scope of the executed operational plan. Throughput/$N_{\mathrm{EFC}}$ are utilization proxies, not a lifetime or cost model. We do not invent an ageing study.

\textbf{(c) Action.} Explicit refusal in metrics and conclusion.

\subsection*{R4.C9 Richer converter/UPS/DC-bus dynamics}
\textbf{(a) Concern.} Verify conclusions with more realistic dynamics.

\textbf{(b) Response.} Only $\tau_b$/SOC on the aggregated plant. Running the same model on the FPGA does not answer this objection.

\textbf{(c) Action.} Model limitations retained.

\subsection*{R4.C10 Statistical robustness}
\textbf{(a) Concern.} More than isolated deterministic scenarios.

\textbf{(b) Response.} 200 pairs and extra windows, with overlapping windows and a single MC physical trajectory disclosed.

\textbf{(c) Action.} MC caveats.

\subsection*{R4.C11 Data-center specificity}
\textbf{(a) Concern.} Show a data-center contribution, not a generic microgrid.

\textbf{(b) Response.} Motivation, mapped AI traces, and 100\%-critical IT load are data-center oriented. The electrical plant remains a generic aggregated source--load--BESS model; multi-UPS and job-class shedding are not modelled.

\textbf{(c) Action.} Related work + critical-load hypothesis.

\subsection*{R4.C12 Qualified conclusions}
\textbf{(a) Concern.} No general superiority or universal zero.

\textbf{(b) Response.} Filtered zero is Group~A; Group~B bound is 350\,kW; long outage: MPC slightly worse on energy; myopic can beat heuristic only after $\Delta u$ retuning.

\textbf{(c) Action.} Conclusion rewrite.

%% =================================================================
\section*{Reviewer 5}

\subsection*{R5.1 Timing precision needed}
\textbf{(a) Concern.} Expand timing analysis and state the accuracy required.

\textbf{(b) Response.} For beating the heuristic peak, both V5 and stochastic pass all 41 points. For the descriptive criterion $P_{\mathrm{un}}^{\max}\le 351$\,kW, V5 passes on $[-20,+4]$\,s. That window is not a pre-registered requirement and is not a universal accuracy spec.

\textbf{(c) Action.} Table~\texttt{tab:err\_time} and the $[-20,+4]$\,s sentence.

\subsection*{R5.2 Why the myopic loses}
\textbf{(a) Concern.} Explain the myopic being worse; retune.

\textbf{(b) Response.} On load bands, myopic $102.78$ vs heuristic $90.98$\,kW at $\Delta u_{\max}=120$. Raising $\Delta u_{\max}$ to 400\,kW/s drives that isolated peak to $\approx 0$. Horizon is therefore not the only cause; $\Delta u$ coefficients also differ.

\textbf{(c) Action.} Fairness paragraph.

\subsection*{R5.3 Align title and all sections with limited scope}
\textbf{(a) Concern.} Title and every section must match the limited event/model scope.

\textbf{(b) Response.} Title now states an aggregated load and detectable events. Abstract/conclusions drop UPS-fidelity and random-outage coverage.

\textbf{(c) Action.} Title, abstract, keywords, conclusion.

\subsection*{R5.4 Quantify SOC, recovery, peak cut, generator coordination}
\textbf{(a) Concern.} Quantify SOC, recovery, peak reduction, backup-generator coordination and nominal benefit.

\textbf{(b) Response.} Peak/SOC/energy numbers are in the tables. Closed-loop generator coordination is \emph{not} simulated; discussion only (R1.8/R1.11).

\textbf{(c) Action.} Tables + discussion-only generator sentence.

\subsection*{R5.5 Unidentified parameters, noise, nonlinearities, run-to-run variability}
\textbf{(a) Concern.} Limits from unidentified parameters, measurement noise, nonlinearities and execution variability.

\textbf{(b) Response.} $\tau_b$ and $\eta$ were not identified on a field battery. Measurement noise is not injected. Solver variability on the i7 is reported as cycle percentiles; FPGA ADMM timing for 1200 steps is not.

\textbf{(c) Action.} Conclusion limitations; no deployment-readiness claim.

\end{document}
